\scanpage{205}\mspage{- 67 -}
\hfill IX

En effet, la dernière hypothèse faite nous ramène au cas où
$G'=A$, $u=\text{identité}$, et où $v_{\handwritten{\eta}}:G_{\handwritten{\eta}}\handwritten{\longrightarrow}A_{\handwritten{\eta}}$ est donc une projection, qu'on se propose de prolonger en une projection $v:G\handwritten{\longrightarrow}A$. Or $G$ est \uline{lisse} comme extension de deux schémas en groupes lisses, de sorte que la conclusion résulte de la définition de la propriété néronienne de $A$.

\uline{Corollaire} 5.10. \uline{Soient} $A_K$ \uline{un schéma abélien sur le corps des fonctions} $K$ \uline{d'un trait} $S$, $\handwritten{\ell}$ \uline{un nombre premier}. \uline{Lorsque} $\handwritten{\ell}=p$, \uline{on suppose} $\mathrm{car.}K=0$,
\uline{pour pouvoir disposer du théorème de TATE déjà invoqué dans} 5.8. \uline{Pour que} $A_K$ \uline{ait bonne réduction sur} $S$, \uline{il faut et il suffit que le pro}-$\handwritten{\ell}$-\uline{groupe de Barsotti-Tate} $T_{\handwritten{\ell}}(A_K)$ ``\uline{ait bonne réduction sur} $S$'', \uline{i}.\uline{e}. \uline{se prolong}e \uline{en un groupe de Ba\deleted{\unsure{t|r}}\unsure{r|t}sotti-Tate sur} $S$.

(C'est le résultat conjecturé dans [8, 4.9].)

La nécessité étant claire, il reste à prouver la suffisance ; notons d'ailleurs que le cas $\handwritten{\ell}\ne p$ a déjà été traité (2.2.9). On peut supposer $S$ hensélien (2.2.9). Supposons d'abord que $A_K$ ait réduction semi-stable sur $S$. Alors 5.8 appliqué à $F=X$ nous montre qu'on a
$T_{\handwritten{\ell}}(A_K)=T_{\handwritten{\ell}}(A_K)^{\mathrm f}$, ce qui signifie aussi que le tore maximal de $A_{\unsure{o|0}}$ est
nul, i.e. que $A_o^o$ est un schéma abélien, de sorte que la conclusion résulte de (2.2.9 (iii$^{\mathrm o}$)). Traitons maintenant le cas général. En vertu de 2.2.9 on peut supposer évidemment $S$ strictement local.
Soir $K'$ une extension galoisienne finie telle que $A_{K'}$ ait réduction semi-stable sur le normalisé $S'$ de $S$ dans $K'$ (3.6). Si $A^*$ désigne le modèle de Néron correspondant, $G=\mathrm{Gal}(K'/K)$ opère sur la situation $(A^*,S')$, donc sur
% archive sheet number: 204
% pencil number at lower centre: 385
% stamp: numbered label 385
